\chapter{Plugins}\label{kap:plugins}

\produkt{} ist zunächst ein CAD-Programm: Sie zeichnen, modellieren, importieren und
exportieren Geometrie. Die eigentlichen Berechnungen – Strahlverfolgung, Lichttransport,
Beugung an Gittern, Scheinwerferauslegung – stecken in eigenständigen
\emph{Rechenkernen}\index{Rechenkern}. Ein \emph{Plugin}\index{Plugin} verbindet einen solchen
Kern mit der Anwendung: Es liest das CAD-Modell, nimmt Ihre Parameter entgegen, startet die
Rechnung und gibt die Ergebnisse so zurück, dass Sie sie direkt am Modell sehen – Strahlen als
Linien, Bestrahlungsstärken als Farbe auf Flächen, Kurven in einem Diagramm, Bilder in einem
eigenen Fenster.

Dieses Kapitel beschreibt alles, was für \emph{alle} Plugins gleich ist: das Plugin-Dock mit
seinen Eingabefeldern, das Starten und Abbrechen von Rechnungen, die Darstellung der
Ergebnisse, das Rückgängigmachen eines Plugin-Laufs und die Einrichtung der
Python-Umgebung. Die einzelnen mitgelieferten Plugins haben eigene Kapitel (Überblick in
Abschnitt~\ref{sec:plugins-mitgeliefert}); den Node-Workflow, in dem mehrere Plugins
zusammenarbeiten, beschreibt Kapitel~\ref{kap:workflow}. Wie Plugins installiert, signiert und
aktualisiert werden, steht in Anhang~\ref{kap:plugin-verwaltung}.

% ---------------------------------------------------------------------------
\section{Was ist ein Plugin?}\label{sec:plugins-konzept}

Ein Plugin ist ein Ordner mit einer Beschreibungsdatei \datei{plugin.json} und Python-Code.
Die Anwendung startet jedes geladene Plugin in einem \emph{eigenen Python-Prozess}
\index{Plugin-Prozess} und spricht mit ihm über eine lokale Verbindung. Das hat für Sie
mehrere Vorteile:

\begin{itemize}
  \item Stürzt eine Rechnung ab oder hängt sie, bleibt \produkt{} bedienbar; nur die Aktion
        des betroffenen Plugins endet mit einer Fehlermeldung.
  \item Eine laufende Rechnung lässt sich immer abbrechen (Abschnitt~\ref{sec:plugin-jobs}).
  \item Jedes Plugin nutzt die Python-Pakete Ihrer Umgebung (z.\,B.\ \code{numpy},
        \code{raytatouille}); die Anwendung selbst bringt kein Python mit
        (Abschnitt~\ref{sec:plugin-einstellungen}).
  \item Große Datenmengen (Millionen Strahlen, Bilder) werden nicht über die Verbindung
        kopiert, sondern als temporäre Dateien übergeben.
\end{itemize}

Ein Plugin kann der Anwendung Folgendes hinzufügen:

\begin{description}
  \item[Aktionen] Befehle wie \gui{Trace rays} oder \gui{Camera image}. Sie erscheinen als
        Knöpfe im Plugin-Dock, im Menü \menu{Plugins} unter dem Namen des Plugins und als
        Werkzeuge in einer eigenen Gruppe der Werkzeugspalte (links).
  \item[Parameter] Eingabefelder im Plugin-Dock (Abschnitt~\ref{sec:plugin-dock}). Die Werte
        werden \emph{mit dem Dokument} gespeichert.
  \item[Ergebnisse] Strahlen, Felder auf Flächen, Kurven, Bilder, Tabellen und
        Balkendiagramme (Abschnitt~\ref{sec:plugin-ergebnisse}).
  \item[Geometrie] Neue oder geänderte Objekte im Modell, z.\,B.\ die Flächen eines geladenen
        optischen Systems. Alle Änderungen eines Laufs bilden einen einzigen Rückgängig-Schritt
        (Abschnitt~\ref{sec:plugin-undo}).
  \item[Blöcke für den Node-Workflow] Bausteine, die im Fenster \gui{Node workflow}
        verschaltet werden (Kapitel~\ref{kap:workflow}).
  \item[Aktionen an Blöcken] Befehle, die nur im Kontextmenü bestimmter Workflow-Blöcke
        stehen, z.\,B.\ \gui{Section $\to$ group} des Raytracers.
\end{description}

\begin{hinweis}
Ein Plugin läuft mit Ihren Benutzerrechten. Die Signaturprüfung
(Anhang~\ref{kap:plugin-verwaltung}) stellt sicher, \emph{von wem} ein Plugin stammt und dass
es unverändert ist – sie ist keine Sandbox. Installieren Sie nur Plugins aus vertrauenswürdiger
Quelle.
\end{hinweis}

\subsection{Start der Plugins}

Beim Programmstart und nach jedem Neuladen (\menu{Plugins,Reload plugins}) sucht die Anwendung
die Plugins in den Plugin-Ordnern (Anhang~\ref{sec:verwaltung-ordner}) und startet jedes
gefundene Plugin sofort. Erst wenn sich das Plugin gemeldet hat, kennt die Anwendung seine
Aktionen und Parameter. Bis dahin steht im Untermenü des Plugins und im Plugin-Dock
\gui{(starting …)} bzw.\ \gui{starting …}. Lässt sich das Plugin nicht starten – meist fehlt
ein Python-Paket oder der Python-Interpreter ist falsch eingestellt –, steht dort
\gui{(not running – see the plugin log)}, und die Ursache steht im Dock \gui{Plugin log}
(Abschnitt~\ref{sec:plugin-log}).

% ---------------------------------------------------------------------------
\section{Das Menü \enquote{Plugins}}\label{sec:plugins-menue}

Das Menü \menu{Plugins} steht in der Menüleiste vor \menu{Help}. Es enthält:

\begin{referenz}{@{}lX@{}}
\toprule
\textbf{Eintrag} & \textbf{Bedeutung} \\
\midrule
\endhead
\gui{Settings…} & Dialog \gui{Plugin settings}: Python-Interpreter, Plugin-Ordner,
  Umgebung prüfen oder anlegen, Entwicklermodus (Abschnitt~\ref{sec:plugin-einstellungen}). \\
\gui{Install plugin…} & Ein Plugin-Paket (\datei{.nfplugin}) installieren
  (Anhang~\ref{sec:verwaltung-installieren}). \\
\gui{Uninstall plugin} & Untermenü mit allen geladenen Plugins (Name und Version). Nur
  installierte Plugins lassen sich hier entfernen; Plugins aus einem Plugin-Ordner sind
  ausgegraut mit dem Zusatz \gui{– from a plugin folder}, der Tooltip nennt den Ordner. \\
\gui{Reload plugins} & Alle Plugins beenden und neu laden, z.\,B.\ nach dem Installieren
  eines Python-Pakets. \\
\gui{Show in the side panels} & Untermenü mit einem Häkchen je Plugin: welche Plugins in der
  Werkzeugspalte und im Plugin-Dock erscheinen (Abschnitt~\ref{sec:plugin-seitenleisten}). \\
\textit{Name des Plugins} & Je Plugin ein Untermenü mit seinen Aktionen; der Tooltip zeigt
  die Beschreibung des Plugins. \\
\bottomrule
\end{referenz}

Sind keine Plugins geladen, steht im Menü der ausgegraute Eintrag
\gui{(no plugins loaded – see Settings)}.

Plugins, die im Entwicklermodus ohne gültige Signatur geladen wurden, erscheinen überall mit
einem Zusatz in eckigen Klammern, z.\,B.\ \gui{Raytracer [unsigned]} oder
\gui{… [signature invalid]} (Anhang~\ref{sec:verwaltung-signaturen}).

% ---------------------------------------------------------------------------
\section{Das Plugin-Dock}\label{sec:plugin-dock}\index{Plugin-Dock}

Das Dock \gui{Plugin} liegt rechts als Reiter neben \gui{Results} (und den Reitern
\gui{Layers} und \gui{Model}); ein- und ausblenden können Sie es im Menü \menu{View}. Es
zeigt immer \emph{ein} Plugin:

\begin{enumerate}
  \item Oben eine Auswahlliste mit allen Plugins, die in den Seitenleisten gezeigt werden.
  \item Darunter die Beschreibung des gewählten Plugins (grau).
  \item Das \emph{Parameterformular}: Eingabefelder, oft in beschriftete Gruppen
        (Rahmen) aufgeteilt. Fährt der Mauszeiger über ein Feld, erklärt ein Tooltip den
        Parameter.
  \item Unten die \emph{Aktionsknöpfe} des Plugins und – sofern es Parameter hat – der flache
        Knopf \gui{Reset parameters}.
\end{enumerate}

\abbildung[0.95\linewidth]{raytracer.png}{Hauptfenster mit dem Raytracer-Plugin: rechts unten
das Plugin-Dock mit Auswahlliste, Beschreibung, Parametern und Aktionsknöpfen; in allen vier
Ansichten die gezeichneten Strahlen und die Farbskala der Bestrahlungsstärke; links die
Werkzeuggruppe des Plugins (Aufnahme im Entwicklermodus, daher \enquote{[unsigned]} und das
rote Banner)}{fig:plugins-dock}

Ist kein Plugin geladen, steht im Dock \gui{No plugins – see Plugins $\to$ Settings…}. Haben Sie
alle Plugins aus den Seitenleisten genommen, steht dort
\gui{All plugins are hidden – see Plugins $\to$ Show in the side panels}.

\subsection{Parameterwerte und Dokument}

Die Werte des Formulars gehören zum \emph{Dokument}: Sie werden je Plugin in der
\datei{.rttf}-Datei gespeichert (Abschnitt \code{plugins}, siehe
Anhang~\ref{kap:dateiformat}) und beim Öffnen wieder eingesetzt. So rechnet ein Dokument nach
dem Öffnen mit genau den Einstellungen, mit denen es gespeichert wurde. Jede Änderung eines
Werts markiert das Dokument als geändert, ist aber \emph{kein} Rückgängig-Schritt.

\gui{Reset parameters} setzt alle Parameter des gezeigten Plugins auf die Standardwerte des
Plugins zurück. Manche Aktionen setzen selbst Parameter, z.\,B.\ wenn ein Plugin eine
Konfigurationsdatei lädt; das Formular zeigt die neuen Werte dann sofort.

Ist das Dock schmal, rückt jedes Eingabefeld unter seine Beschriftung, damit nichts
abgeschnitten wird.

\subsection{Die Feldtypen}\label{sec:plugin-parameter}\index{Parameterformular}

Welche Felder erscheinen, legt das Plugin fest. Die Anwendung kennt die folgenden Feldtypen:

\begin{referenz}{@{}p{0.2\linewidth}p{0.3\linewidth}X@{}}
\toprule
\textbf{Typ} & \textbf{Aussehen} & \textbf{Bedienung} \\
\midrule
\endhead
Zahl & Drehfeld mit Einheit (z.\,B.\ \code{0.0001000 mm}) &
  Wert tippen oder mit den Pfeilen ändern. Dezimaltrennzeichen ist immer der \emph{Punkt}
  (wie in der Befehlszeile), unabhängig von der Spracheinstellung des Systems. Das Plugin kann
  Grenzen, Schrittweite und Nachkommastellen vorgeben (Standard: vier Nachkommastellen). \\
Ganzzahl & Drehfeld, ggf.\ mit Einheit & wie Zahl, ohne Nachkommastellen. \\
Ja/Nein & Kontrollkästchen & anklicken. \\
Auswahl & Auswahlliste & einen der vorgegebenen Einträge wählen. \\
Text & einzeiliges Textfeld & frei tippen. \\
Datei & Textfeld mit Knopf \gui{…} & Pfad tippen oder mit \gui{…} (Tooltip
  \gui{Choose a file}) im Dateidialog wählen. Ein leeres Feld zeigt grau
  \gui{choose a file with …}. Bei Ausgabedateien öffnet \gui{…} einen Speichern-Dialog. \\
Dateiliste & Liste mit \gui{Add…} und \gui{Remove} & \gui{Add…} öffnet den Dateidialog,
  in dem Sie \emph{mehrere} Dateien auf einmal wählen können; eine Datei, die schon in der
  Liste steht, wird nicht doppelt aufgenommen. Die Liste zeigt die Dateinamen, der Tooltip
  den vollen Pfad. \gui{Remove} entfernt die markierten Einträge (Mehrfachauswahl mit
  \keys{Shift} oder \keys{Strg}). \\
Zahlenliste & Textfeld & Zahlen durch Kommas getrennt, Dezimalpunkt, z.\,B.\
  \code{486.1, 587.6, 656.3}. Eine ungültige Eingabe bekommt einen roten Rahmen. \\
Punkt & drei Felder \code{x}, \code{y}, \code{z} und Knopf \gui{Pick} & Koordinaten tippen
  (drei Nachkommastellen) oder mit \gui{Pick} in einer Ansicht klicken
  (Abschnitt~\ref{sec:plugin-punkt}). \\
Objekt & Name(n), Knöpfe \gui{Use selection} und \gui{$\times$} & übernimmt Objekte aus
  der aktuellen Auswahl (Abschnitt~\ref{sec:plugin-objekt}). \\
Liste von Einträgen & Liste mit \gui{Add}, \gui{Edit…}, \gui{Remove}, \gui{$\uparrow$},
  \gui{$\downarrow$} & geordnete Einträge, jeder ein kleines Formular, z.\,B.\ die Schichten
  eines Gitters (Abschnitt~\ref{sec:plugin-liste}). \\
Tabelle & Tabelle mit \gui{Add row} und \gui{Remove row} & direkt in den Zellen bearbeiten
  (Abschnitt~\ref{sec:plugin-tabelle}). \\
JSON & mehrzeiliges Textfeld & beliebige Daten im JSON-Format; ungültiger Text bekommt
  einen roten Rahmen. \\
\bottomrule
\end{referenz}

\begin{warnung}
Solange ein Feld rot markiert ist (Zahlenliste, Tabellenzelle oder JSON-Text ungültig), tut
ein Klick auf einen Aktionsknopf des Plugins \emph{nichts}. Korrigieren Sie zuerst die rote
Eingabe. Erst gültige Eingaben werden übernommen; bis dahin gilt der letzte gültige Wert.
\end{warnung}

\subsubsection{Einen Punkt picken}\label{sec:plugin-punkt}

Punkt-Parameter (z.\,B.\ ein Ursprung oder Zielpunkt) können Sie direkt im Modell abgreifen:

\begin{schritte}
  \item Klicken Sie neben dem Punkt-Parameter auf \gui{Pick}. Die Befehlszeile zeigt
        \gui{Pick a point for the plugin parameter (snaps active, Esc cancels)}.
  \item Klicken Sie in einer beliebigen Ansicht auf den gewünschten Punkt. Die Objektfänge
        sind aktiv wie beim Zeichnen (Kapitel~\ref{kap:modellieren}): Endpunkte, Mittelpunkte,
        Zentren usw.\ werden gefangen, \keys{Tab} wechselt zwischen Fangkandidaten.
        Statt zu klicken, können Sie die Koordinaten auch in die Befehlszeile tippen.
  \item Die Felder \code{x}, \code{y}, \code{z} übernehmen den Punkt.
\end{schritte}

Mit \keys{Esc} brechen Sie das Picken ab; die alten Werte bleiben.

\subsubsection{Objekte aus der Auswahl übernehmen}\label{sec:plugin-objekt}

Objekt-Parameter verweisen auf Objekte des Modells, etwa einen Detektor oder die Körper einer
Gitterschicht. Das Plugin legt fest, welche Objektarten erlaubt sind (Punkt, Kurve, Fläche,
getrimmte Fläche, Körper/Shell) und ob ein oder mehrere Objekte gemeint sind.

\begin{schritte}
  \item Wählen Sie die Objekte in einer Ansicht oder im Modellbaum aus. Teilelemente (einzelne
        Flächen, Kanten, Eckpunkte eines Körpers) zählen nicht.
  \item Klicken Sie auf \gui{Use selection}. Das Feld zeigt nun die Namen der übernommenen
        Objekte (ohne Namen: \code{\#}\textit{Nummer}). Bei einem Parameter für ein einzelnes
        Objekt wird das erste passende Objekt der Auswahl genommen.
\end{schritte}

Passt kein Objekt der Auswahl, meldet die Statuszeile
\gui{Select suitable objects first (…)} mit den erlaubten Arten. \gui{$\times$} leert das
Feld (Anzeige \gui{---}). Wird ein übernommenes Objekt später gelöscht, zeigt das Feld
\code{\#}\textit{Nummer}~\gui{(missing)}.

\subsubsection{Listen von Einträgen}\label{sec:plugin-liste}

Strukturierte Parameter – etwa der Schichtaufbau eines Gitters, Materialien oder die Terme einer
Zielfunktion – erscheinen als Liste. Jede Zeile ist ein Eintrag, beschriftet mit seiner Nummer
und einem kennzeichnenden Feld, z.\,B.\ \code{1. homogeneous}.

\begin{itemize}
  \item \gui{Add} öffnet den Dialog \gui{Edit entry} mit den Standardwerten eines neuen
        Eintrags; \gui{OK} hängt ihn am Ende an.
  \item \gui{Edit…} (oder ein Doppelklick auf den Eintrag) öffnet den markierten Eintrag im
        selben Dialog. Der Dialog ist wieder ein Formular mit allen Feldtypen – auch Listen
        und Tabellen in Listen sind möglich.
  \item \gui{Remove} löscht den markierten Eintrag.
  \item \gui{$\uparrow$} und \gui{$\downarrow$} verschieben ihn um eine Position.
\end{itemize}

Der Dialog lässt sich mit \gui{OK} erst schließen, wenn kein Feld darin rot markiert ist.

\subsubsection{Tabellen}\label{sec:plugin-tabelle}

Tabellen-Parameter (z.\,B.\ die Lamellen einer Gitterschicht oder die Elemente einer
Streuscheibe) bearbeiten Sie direkt in den Zellen. Die Spaltenköpfe nennen Größe und Einheit.
\gui{Add row} hängt eine Zeile mit Standardwerten an, \gui{Remove row} löscht die Zeile mit
dem Cursor. Eine Zelle, die in einer Zahlenspalte keine (ganze) Zahl enthält oder in einer
Auswahlspalte keinen erlaubten Wert, wird rot hinterlegt und blockiert die Aktionen.

% ---------------------------------------------------------------------------
\section{Aktionen und Jobs}\label{sec:plugin-jobs}\index{Job}

Eine \emph{Aktion} ist ein Befehl eines Plugins. Sie starten sie auf drei gleichwertigen
Wegen:

\begin{itemize}
  \item mit dem Knopf im Plugin-Dock,
  \item im Menü \menu{Plugins,\textit{Plugin},\textit{Aktion}},
  \item mit dem Werkzeug in der Gruppe des Plugins in der Werkzeugspalte (der Tooltip nennt
        die Aktion).
\end{itemize}

Die Aktion erhält die aktuellen Parameterwerte und die aktuelle Auswahl. Während sie läuft,
spricht man von einem \emph{Job}:

\begin{itemize}
  \item Die Statuszeile zeigt rechts einen Fortschrittsbalken und den Knopf \gui{Cancel}, links
        den Text \textit{Plugin}\code{:} \textit{Aktion}~\code{…} und danach die Meldungen, die
        das Plugin während der Rechnung schickt.
  \item Die Aktionsknöpfe im Plugin-Dock sind gesperrt (\gui{Reset parameters} bleibt
        bedienbar).
  \item Es läuft immer nur \emph{ein} Plugin-Job. Starten Sie einen zweiten, meldet die
        Statuszeile \gui{A plugin job is still running}.
  \item Die Anwendung bleibt bedienbar: Sie können drehen, zoomen, Ergebnisse ein- und
        ausblenden. Ändert eine Aktion das Modell laufend (z.\,B.\ eine Optimierung), sehen Sie
        die Änderungen sofort.
\end{itemize}

\subsection{Abbrechen}\label{sec:plugin-abbrechen}\index{Abbrechen}

Sie brechen einen Job mit \gui{Cancel} in der Statuszeile oder mit \keys{Esc} ab. Die Taste
\keys{Esc} (\menu{Edit,Select None / Cancel}) wirkt in dieser Reihenfolge: Läuft ein
Zeichenbefehl, wird er abgebrochen; sonst, wenn ein Plugin-Job läuft, dieser; sonst wird die
Auswahl aufgehoben.

Das Plugin erhält die Bitte, aufzuhören, und beendet seine Rechnung an der nächsten geeigneten
Stelle (z.\,B.\ nach dem nächsten Iterationsschritt). Reagiert es nicht innerhalb von
\emph{5~Sekunden}, beendet die Anwendung den Plugin-Prozess; das Log meldet dann
\gui{the plugin did not stop – terminating its process}. Das Plugin wird bei der nächsten
Aktion neu gestartet.

\begin{hinweis}
Auch ein abgebrochener Lauf bildet einen Rückgängig-Schritt: Was bis zum Abbruch am Modell
geändert wurde, bleibt erhalten und lässt sich mit \keys{Strg+Z} zurücknehmen
(Abschnitt~\ref{sec:plugin-undo}).
\end{hinweis}

\subsection{Ende eines Jobs}

Nach dem Lauf meldet die Statuszeile:

\begin{referenz}{@{}p{0.42\linewidth}X@{}}
\toprule
\textbf{Meldung} & \textbf{Bedeutung} \\
\midrule
\endhead
\textit{Plugin}\code{: done} oder \textit{Plugin}\code{:} \textit{Text} & Erfolg; den Text
  liefert das Plugin, z.\,B.\ eine Zusammenfassung. Die Meldung verschwindet nach acht
  Sekunden. \\
\textit{Plugin}\code{: cancelled} & Der Job wurde abgebrochen. \\
\textit{Plugin}\code{ failed:} \textit{Grund} & Fehler. Das Dock \gui{Plugin log} öffnet
  sich mit den Einzelheiten (bei Programmfehlern mit dem Python-Traceback). \\
\bottomrule
\end{referenz}

Fehler, die Sie selbst beheben können – keine Systemdatei gewählt, eine fehlende Eingabe,
eine nicht lesbare Datei –, zeigt die Anwendung stattdessen als Meldungsfenster mit einem
verständlichen Text, ohne Traceback.

Hat ein Lauf in ein \emph{leeres} Dokument Objekte eingefügt (z.\,B.\ das erste Laden eines
optischen Systems), zoomt die Anwendung anschließend alle Ansichten auf das Ergebnis.

\subsection{Das Plugin-Log}\label{sec:plugin-log}\index{Plugin-Log}

Das Dock \gui{Plugin log} liegt unten als Reiter neben \gui{Plot}; zu Beginn ist es
ausgeblendet (\menu{View,Plugin log}). Es sammelt alle Meldungen der Plugins, jede Zeile mit
der Plugin-Kennung in eckigen Klammern, z.\,B.\ \code{[raytracer] …}. Dort finden Sie:

\begin{itemize}
  \item ausführliche Protokolle der Rechnungen (Anzahl Strahlen, verwendete Verfahren,
        Rauschen, Warnungen),
  \item Fehlermeldungen mit Traceback,
  \item Startprobleme eines Plugins, z.\,B.\ ein fehlendes Python-Paket,
  \item die Meldung, wenn ein Plugin-Prozess unerwartet endet:
        \code{plugin '}\textit{Name}\code{' stopped (exit code} \textit{n}\code{)}
        mit den letzten Ausgaben des Prozesses.
\end{itemize}

Das Log behält die letzten 5000 Zeilen.

% ---------------------------------------------------------------------------
\section{Ergebnisse}\label{sec:plugin-ergebnisse}\index{Ergebnis}

Jeder Lauf einer Aktion kann einen \emph{Ergebnissatz}\index{Ergebnissatz} (engl.\ run) anlegen,
z.\,B.\ \enquote{Trace 1}, und darin beliebig viele Ergebnisse ablegen. Es gibt sechs Arten:

\begin{referenz}{@{}p{0.2\linewidth}p{0.17\linewidth}X@{}}
\toprule
\textbf{Art} & \textbf{Typ-Spalte} & \textbf{Darstellung} \\
\midrule
\endhead
Strahlen / Linien & \gui{Rays / lines} & Viele farbige Polylinien in allen Ansichten, z.\,B.\
  Strahlen, eingefärbt nach Wellenlänge, Feld oder Pfad. \\
Feld auf Fläche & \gui{Field on face} & Eine Werteverteilung (oder ein Farbbild) als Textur auf
  einer Fläche des Modells, z.\,B.\ die Bestrahlungsstärke auf einem Detektor. \\
Kurve & \gui{Curve} & x/y-Kurve im Dock \gui{Plot}, z.\,B.\ der Verlauf der Zielfunktion
  einer Optimierung; kann während der Rechnung wachsen. \\
Bild & \gui{Image} & 2D-Bild mit physikalischen Achsen im Fenster \gui{Image}, z.\,B.\ ein
  Kamerabild. \\
Tabelle & \gui{Table} & Tabelle im Dock \gui{Table}, z.\,B.\ die Beugungseffizienzen je
  Ordnung. \\
Balken & \gui{Bars} & Balkendiagramm im Dock \gui{Plot}. \\
\bottomrule
\end{referenz}

\subsection{Das Results-Dock}\label{sec:results-dock}\index{Results-Dock}

Das Dock \gui{Results} (rechts, Reiter neben \gui{Model}) listet alle Ergebnissätze als Baum
mit den Spalten \gui{Result} und \gui{Type}. Ein Ergebnissatz zeigt in der Spalte \gui{Type}
die Kennung des Plugins, das ihn angelegt hat; darunter stehen seine Ergebnisse mit ihrer Art.

\begin{itemize}
  \item \textbf{Ein- und ausblenden:} Das Kontrollkästchen vor einem Ergebnissatz blendet alle
        seine Ergebnisse ein oder aus, das Kästchen vor einem Ergebnis nur dieses. So lassen sich
        z.\,B.\ die Strahlen eines Arms eines Interferometers ausblenden.
  \item \textbf{Löschen:} Rechtsklick auf einen Ergebnissatz $\to$ \gui{Delete run}.
  \item \textbf{Speichern:} Rechtsklick auf einen Ergebnissatz $\to$ \gui{Save run…} oder auf
        ein Ergebnis $\to$ \gui{Save…}; dasselbe macht die Schaltfläche \gui{Save…} unter dem
        Baum für die gewählte Zeile (Abschnitt~\ref{sec:ergebnisse-speichern}).
  \item \textbf{Strahlen in CAD-Kurven umwandeln:} Rechtsklick auf ein Linien-Ergebnis
        $\to$ \gui{Convert to CAD curves…} (Abschnitt~\ref{sec:strahlen-kurven}).
  \item Der Tooltip eines Linien-Ergebnisses nennt die Zahl der Polylinien.
\end{itemize}

\begin{hinweis}
Ergebnisse sind Anzeigedaten und gehören nicht zum Modell: Sie werden nicht gefangen, nicht
ausgewählt und – mit einer Ausnahme – nicht mit dem Dokument gespeichert (wohl aber auf Wunsch als Dateien,
Abschnitt~\ref{sec:ergebnisse-speichern}). Die Ausnahme sind die
Ergebnisse der Anzeigeblöcke des Node-Workflows; sie stehen nach dem Öffnen im Ergebnissatz
\gui{Saved results} (Abschnitt~\ref{sec:workflow-gespeichert}).
\end{hinweis}

\subsection{Strahlen in den Ansichten}

Linien-Ergebnisse werden standardmäßig \emph{über} dem Modell gezeichnet, also ohne
Verdeckung – Strahlen im Inneren einer Linse bleiben so sichtbar. Ein Plugin kann für ein
Linien-Ergebnis auch die normale Verdeckung wählen. Ein verlorener Strahl endet dort, wo er das
System verlässt oder abgeschattet wird. Trägt ein Linien-Ergebnis Werte (z.\,B.\ die
Wellenlänge), wird es über eine Farbskala eingefärbt, sonst einfarbig.

\subsubsection{Strahlen als CAD-Kurven}\label{sec:strahlen-kurven}

Strahlen sind Anzeigedaten; zum Messen oder Konstruieren brauchen Sie echte Kurven:

\begin{schritte}
  \item Rechtsklick auf das Linien-Ergebnis im Dock \gui{Results} $\to$
        \gui{Convert to CAD curves…}.
  \item Der Dialog \gui{Convert to curves} fragt \gui{Every n-th ray (of N):}. Vorgeschlagen
        wird ein Wert, der etwa 50 Kurven ergibt.
  \item Nach \gui{OK} entstehen Polylinien mit den Namen \code{Ray 0}, \code{Ray 5}, … –
        höchstens 2000 Stück. Ein verlorener Strahl endet an seinem letzten gültigen Punkt.
\end{schritte}

Die neuen Kurven sind normale CAD-Objekte (fangbar, messbar, exportierbar) und bilden einen
Rückgängig-Schritt \enquote{Rays to curves}.

\subsubsection{Zeichnungen, die sich ersetzen}

Ein Plugin kann einem Linien-Ergebnis einen \emph{Schlüssel} geben. Ein neues Linien-Ergebnis
mit demselben Schlüssel ersetzt dann das frühere: Dessen Linien verschwinden, und ein
Ergebnissatz, der dadurch leer wird, verschwindet mit. So zeichnet \gui{Show camera} des Path
Tracers nach einer Änderung der Kameraparameter die Kamera neu, statt eine zweite darüber zu legen.

\subsection{Ergebnisse speichern}\label{sec:ergebnisse-speichern}\index{Ergebnis!speichern}\index{CSV}\index{PNG}

Ergebnisse lassen sich als Dateien speichern -- Bilder als PNG, Zahlen als CSV (Komma als
Trennzeichen, Punkt als Dezimalzeichen, UTF-8, Gleitkommazahlen mit neun signifikanten Stellen,
\code{nan} für ungültige Werte):

\begin{referenz}{@{}p{0.2\linewidth}X@{}}
\toprule
\textbf{Art} & \textbf{Dateien} \\
\midrule
\endhead
Bild & PNG: ein Bildpunkt je Wert, mit der Farbskala (lineare Skala; im Image-Fenster wie
  angezeigt, auch logarithmisch), ein Farbbild in seinen Farben. CSV: Kommentarzeilen mit
  \code{\#} (Name, Einheit, Ausdehnung, Achsen), dann eine Zeile je Bildzeile von $y_0$ aufwärts,
  eine Spalte je $x$ -- \code{numpy.loadtxt(datei, delimiter=",")} liest das als Feld
  [Zeile, Spalte]. \\
Feld auf Fläche & PNG mit der Farbskala über $(u, v)$, $v$ nach oben; CSV wie ein Bild, eine
  Zeile je $v$, eine Spalte je $u$. \\
Kurve & CSV: Kopfzeile mit den Achsbeschriftungen, dann \code{x,y} je Zeile. \\
Tabelle & CSV: Spaltenköpfe, dann die Zeilen. \\
Balken & CSV: \code{label,}\textit{y-Beschriftung}, dann ein Balken je Zeile. \\
Strahlen / Linien & CSV: \code{polyline,x,y,z} (und \code{value}, wenn die Punkte Werte tragen),
  ein Punkt je Zeile; \code{nan}, wo ein Strahl verloren ging. \\
\bottomrule
\end{referenz}

\gui{Save…} auf einem Ergebnis fragt nach dem Dateinamen; bei Bildern und Feldern wählen Sie
im Dateityp zwischen PNG und CSV. \gui{Save run…} fragt nach einem Ordner und legt darin einen
Unterordner mit dem Namen des Ergebnissatzes an, mit einer Datei je Ergebnis (Bilder und Felder
als PNG \emph{und} CSV); gleichnamige Ergebnisse werden durchnummeriert (\texttt{orders (2).csv}),
Zeichen, die Dateisysteme nicht erlauben, durch \code{\_} ersetzt. Die Statusleiste meldet, wie
viele Dateien wohin geschrieben wurden.

\subsection{Felder auf Flächen und Farbskalen}\label{sec:felder-farbskalen}\index{Farbskala}

Ein Feld auf einer Fläche ist ein Raster von Werten über den normierten Flächenparametern
$(u, v)$ einer Fläche (bei Körpern einer bestimmten Teilfläche). Die Anwendung färbt die Werte
über eine \emph{Farbskala} und legt das Ergebnis als Textur auf die Fläche. Im Darstellungsmodus
\gui{Shaded} wird die Textur beleuchtet, in \gui{Surfaces} und \gui{Wireframe} flach gezeigt.
Ein Plugin kann statt Werten auch ein fertiges Farbbild auf die Fläche legen.

In jeder Ansicht erscheint unten rechts eine \emph{Legende} mit Name, Einheit, Farbbalken und
Minimum/Maximum (Abbildung~\ref{fig:plugins-dock}). Sie gilt für das erste sichtbare Feld
mit Werten; gibt es keines, für das erste sichtbare, nach Werten gefärbte Linien-Ergebnis.
Blenden Sie Ergebnisse im Dock \gui{Results} aus, um die Legende eines anderen zu sehen.

Die Anwendung kennt fünf Farbskalen; welche verwendet wird, bestimmt das Plugin:

\begin{referenz}{@{}lX@{}}
\toprule
\textbf{Name} & \textbf{Verlauf und typische Verwendung} \\
\midrule
\endhead
\code{viridis} & dunkelviolett $\to$ blau $\to$ grün $\to$ gelb; gleichmäßig wahrgenommene
  Helligkeit, allgemeiner Standard (auch für unbekannte Namen). \\
\code{inferno} & schwarz $\to$ violett $\to$ rot $\to$ hellgelb; Standard für Felder und
  Bilder (Intensitäten, Bestrahlungsstärken). \\
\code{plasma} & dunkelblau $\to$ violett $\to$ orange $\to$ gelb. \\
\code{coolwarm} & blau $\to$ hellgrau $\to$ rot; für Werte mit Vorzeichen (Abweichungen). \\
\code{grey} & schwarz $\to$ weiß (auch \code{gray}). \\
\bottomrule
\end{referenz}

Der Wertebereich der Farbskala ist entweder vom Plugin vorgegeben oder reicht vom kleinsten
zum größten Wert der Daten.

\subsection{Das Plot-Dock}\label{sec:plot-dock}\index{Plot-Dock}

Das Dock \gui{Plot} liegt unten und erscheint von selbst mit der ersten Kurve. Es zeigt alle
sichtbaren Kurven in einem gemeinsamen Diagramm, jede in einer eigenen Farbe, mit einer Legende
(Namen) oben links und Teilstrichen an beiden Achsen; unter der x-Achse steht die Beschriftung
der ersten Kurve. Lücken in einer Kurve (ungültige Werte) werden nicht verbunden. Kurven, die ein
Plugin während der Rechnung verlängert (z.\,B.\ die Zielfunktion je Iteration), wachsen
live.

Ist keine Kurve sichtbar, aber ein Balkendiagramm, zeigt das Dock das neueste Balkendiagramm
mit Titel, y-Beschriftung und den Beschriftungen unter den Balken. Ohne beides steht dort
\gui{No curves (e.g. the merit function of an optimisation)}.

\begin{tipp}
Um nur die Kurven eines Laufs zu sehen, blenden Sie die anderen Ergebnissätze im Dock
\gui{Results} aus.
\end{tipp}

\subsection{Das Table-Dock}\label{sec:table-dock}

Tabellen erscheinen im Dock \gui{Table} (unten), das sich mit der ersten Tabelle öffnet. Gibt
es mehrere sichtbare Tabellen, wählen Sie oben in einer Auswahlliste, welche gezeigt wird; eine
neu hinzugekommene Tabelle wird automatisch gezeigt. Zahlen stehen rechtsbündig; Gleitkommazahlen
schreibt das Plugin mit zehn signifikanten Stellen. Die Spaltenbreiten lassen sich mit der Maus
ändern.

\subsection{Das Image-Fenster}\label{sec:image-fenster}\index{Image-Fenster}

Zweidimensionale Bilder – etwa das Kamerabild des Raytracers oder des Path Tracers – zeigt das
Fenster \gui{Image}. Es öffnet sich \emph{schwebend} in der Mitte des Hauptfensters, sobald ein
neues Bild ankommt. Wie jedes Dock können Sie es an einer Seite andocken, indem Sie seine
Titelleiste dorthin ziehen; \menu{View,Image} blendet es ein und aus.

\begin{description}
  \item[Auswahl] Gibt es mehrere Bilder, wählen Sie oben in der Liste
        (\enquote{\textit{Ergebnissatz} – \textit{Bild}}). Ein neues Bild kommt automatisch
        nach vorn; sonst bleibt Ihre Wahl.
  \item[Achsen] Das Bild hat physikalische Achsen (z.\,B.\ in mm) mit Teilstrichen und
        Beschriftungen; oben steht der Titel. Die Pixel bleiben auch vergrößert scharfkantig
        (keine Glättung).
  \item[Farbbalken] Rechts zeigt ein Farbbalken die Farbskala mit größtem und kleinstem Wert
        und darunter die Einheit.
  \item[log] Das Kontrollkästchen \gui{log} schaltet auf eine logarithmische Farbskala um
        (Tooltip: \gui{Logarithmic colour scale (faint light next to a bright spot)}). Sie
        reicht vom Spitzenwert vier Dekaden nach unten, höchstens bis zum kleinsten positiven
        Wert; der Farbbalken ist dann in Zehnerpotenzen beschriftet (\code{1e-2.0} usw.).
  \item[Ablesen] Unter dem Bild steht die Position des Mauszeigers und der Wert dort,
        z.\,B.\ \code{x 0.0123  y -0.0040   1.52e+03 W/mm$^2$}. Außerhalb des Bildes steht dort
        die Beschreibung des Bildes (z.\,B.\ Zahl der Strahlen).
  \item[Zoomen und Verschieben] Das Mausrad zoomt um den Mauszeiger (bis auf etwa vier Pixel
        Breite), Ziehen mit der linken Maustaste verschiebt den Ausschnitt, ein Doppelklick
        zeigt wieder das ganze Bild.
  \item[Speichern] Die Schaltfläche \gui{Save…} oben rechts speichert das gezeigte Bild:
        \gui{Image as PNG…} -- ein Bildpunkt je Wert, mit Farbskala und \gui{log} wie
        angezeigt; \gui{View with axes as PNG…} -- das Fenster so, wie es aussieht, mit Achsen,
        Titel, Farbbalken und dem aktuellen Ausschnitt; \gui{Values as CSV…} -- die Zahlenwerte
        (Abschnitt~\ref{sec:ergebnisse-speichern}).
\end{description}

\paragraph{Farbbilder.} Ein Plugin kann ein Bild auch in Farbe liefern (z.\,B.\ das
Kamerabild des Path Tracers bei mehreren Wellenlängen). Es wird so gezeigt, wie es ist – ohne
Farbskala, ohne Farbbalken; \gui{log} ist dann ausgegraut. Die Ablesung unter dem Bild zeigt
weiterhin den Zahlenwert, den das Plugin zum Bild mitgeliefert hat.

Ohne Bild zeigt das Fenster \gui{No image – e.g. Raytracer $\to$ Camera image}.

\begin{hinweis}
Die Bilder und Kurven der Anzeigeblöcke des Node-Workflows erscheinen nicht im allgemeinen
Image-Fenster bzw.\ Plot-Dock, sondern in einem eigenen Fenster je Anzeigeblock
(Abschnitt~\ref{sec:workflow-anzeigen}).
\end{hinweis}

\subsection{Ausprobieren ohne Plugin: Demo-Ergebnisse}

Um die Ergebnisdarstellung kennenzulernen, brauchen Sie kein Plugin: Wählen Sie eine Fläche
oder einen Körper und rufen Sie \menu{Help,Demo results on the selected object} auf. Es
entsteht ein Ergebnissatz \gui{Demo run} mit Strahlen (\gui{Rays}, in nm eingefärbt), einem
Feld auf dem gewählten Objekt (\gui{Irradiance}, W/m$^2$) und einer Kurve (\gui{Merit} über
\gui{iteration}). Ohne Auswahl meldet die Statuszeile \gui{Select a surface or solid first}.

% ---------------------------------------------------------------------------
\section{Rückgängig machen eines Plugin-Laufs}\label{sec:plugin-undo}\index{Rueckgaengig@Rückgängig}

Alles, was ein Lauf am \emph{Modell} ändert, bildet \emph{einen} Rückgängig-Schritt, benannt
nach dem Plugin:

\begin{itemize}
  \item neu eingefügte und gelöschte Objekte,
  \item geänderte Geometrie (auch wenn sie während einer Optimierung live viele Male ersetzt
        wurde – zurückgenommen wird auf den Zustand vor dem Lauf),
  \item geänderte Attribute,
  \item neu angelegte und gelöschte Gruppen des Modellbaums,
  \item Änderungen am Node-Workflow, die das Plugin vornimmt (z.\,B.\ \gui{Workflow from system},
        \gui{New Raytatouille system}, \gui{Section $\to$ group} oder das Zurückschreiben einer
        Optimierung in die Blöcke); reine Prüfergebnisse (Diagnosen, Plan) zählen nicht dazu.
\end{itemize}

Das gilt auch für abgebrochene oder fehlgeschlagene Läufe. Mit \menu{Edit,Undo}
(\keys{Strg+Z}) nehmen Sie den ganzen Lauf zurück, mit \menu{Edit,Redo}
(\keys{Strg+Shift+Z}) stellen Sie ihn wieder her. Ein Lauf, der das Modell nicht ändert
(z.\,B.\ eine reine Berechnung mit Ergebnissen), erzeugt keinen Rückgängig-Schritt.

\begin{hinweis}
Nicht rückgängig gemacht werden: Ergebnisse (löschen Sie sie im Dock \gui{Results}) und
Änderungen der Parameterwerte im Plugin-Dock. Manche Plugins merken sich zusätzlich den Stand
ihrer eigenen Daten; der Raytracer etwa rechnet nach einem Rückgängig mit dem System vor der
Optimierung weiter (Kapitel~\ref{kap:raytracer}).
\end{hinweis}

% ---------------------------------------------------------------------------
\section{Plugin-Einstellungen}\label{sec:plugin-einstellungen}\index{Plugin-Einstellungen}

\menu{Plugins,Settings…} öffnet den Dialog \gui{Plugin settings}. Hier legen Sie fest, mit
welchem Python die Plugins laufen und wo sie gesucht werden.

\abbildung[0.8\linewidth]{plugin-settings.png}{Der Dialog \gui{Plugin settings} (Aufnahme aus
einer Entwicklungsumgebung): Interpreter, Umgebungsprüfung, Plugin-Ordner,
Entwicklermodus und Problemliste}{fig:plugins-einstellungen}

\begin{referenz}{@{}p{0.3\linewidth}X@{}}
\toprule
\textbf{Element} & \textbf{Bedeutung} \\
\midrule
\endhead
\gui{Python interpreter:} & Pfad zum Python-Programm, mit dem alle Plugins laufen; \gui{…}
  öffnet einen Dateidialog. Es muss Python 3.10 oder neuer sein. \\
\gui{nfcad package:} & Wo die Anwendung ihr Python-Paket \code{nfcad} gefunden hat (die
  Verbindung zwischen Anwendung und Plugins); \gui{not found} deutet auf eine beschädigte
  Installation. \\
\gui{Check environment} & Prüft den Interpreter im Feld (auch vor \gui{OK}) und die Pakete,
  die die geladenen Plugins brauchen; Ergebnis im Textfeld darunter. \\
\gui{Create environment…} & Legt eine neue virtuelle Umgebung (venv) an und installiert alle
  benötigten Pakete hinein. \\
\gui{Installed plugins (…):} & Der Ordner, in den \gui{Install plugin…} installiert (nur
  Anzeige, Text markierbar). \\
\gui{Further plugin folders} & Weitere Ordner; jeder Unterordner mit einer
  \datei{plugin.json} ist ein Plugin. \gui{Add folder…} fügt einen hinzu, \gui{Remove}
  entfernt den markierten. \\
\gui{Trusted publisher key:} & Der eingebaute öffentliche Schlüssel, dessen Signaturen
  akzeptiert werden (Anhang~\ref{sec:verwaltung-signaturen}). \\
\gui{Developer mode: also load unsigned plugins} & Entwicklermodus
  (Abschnitt~\ref{sec:plugin-entwicklermodus}). \\
\gui{Problems:} & Erscheint nur, wenn es beim Laden Probleme gab (orange): nicht geladene oder
  unsignierte Plugins, doppelte Kennungen, fehlerhafte \datei{plugin.json}. \\
\bottomrule
\end{referenz}

Mit \gui{OK} übernimmt die Anwendung Interpreter, Ordner und Entwicklermodus und lädt alle
Plugins neu. Ein geänderter Interpreter beendet die laufenden Plugin-Prozesse.

\subsection{Welches Python wird verwendet?}\label{sec:python-suche}

Solange Sie keinen Interpreter eingetragen haben, sucht die Anwendung in dieser Reihenfolge:

\begin{enumerate}
  \item den Inhalt der Umgebungsvariablen \code{NFCAD\_PYTHON}
        (Anhang~\ref{kap:programme}),
  \item in einer Entwicklungsversion die \datei{.venv} des Quellcode-Ordners,
  \item unter Windows \code{python}, \code{py} und \code{python3} im Suchpfad – die Platzhalter
        des Microsoft Store (Ordner \datei{WindowsApps}) werden übersprungen, weil sie nur den
        Store öffnen; unter macOS und Linux \code{python3} und \code{python},
  \item zuletzt einfach \code{python3}.
\end{enumerate}

Ein im Dialog eingetragener Interpreter hat immer Vorrang.

\subsection{Umgebung prüfen}\label{sec:umgebung-pruefen}

\gui{Check environment} startet den Interpreter im Feld und versucht, jedes Paket zu
importieren, das die geladenen Plugins in ihrer \datei{plugin.json} verlangen (dazu immer
\code{numpy}). Das Ergebnis sieht etwa so aus:

\begin{lstlisting}[escapeinside={(*}{*)}]
Python 3.11.15  ok
(*\checkmark*) numpy  2.4.6
(*\checkmark*) raytatouille>=0.9  <Version>
(*$\times$*) fouriergrating>=0.18  ModuleNotFoundError: No module named 'fouriergrating'

Install the missing packages with:
"/Users/.../python-env/bin/python" -m pip install "fouriergrating>=0.18"
\end{lstlisting}

Die erste Zeile nennt die Python-Version; ist sie älter als 3.10, steht dort
\code{Python 3.9.6 is too old (3.10 or newer needed)}. Jede weitere Zeile zeigt ein Paket mit
\checkmark{} (vorhanden, mit Version) oder $\times$ (fehlt, mit der Fehlermeldung des Imports
oder \code{version … installed, >=… needed}). Am Ende steht – falls etwas fehlt – der
fertige \code{pip}-Befehl zum Kopieren. Pakete, deren pip-Name vom Importnamen abweicht, nennt
der Befehl mit dem pip-Namen (z.\,B.\ \code{photonpower-fieldtracing} für das Modul
\code{fieldtracing}). Lässt sich der Interpreter nicht starten, steht dort
\code{… cannot be started}.

\subsection{Umgebung anlegen}\label{sec:umgebung-anlegen}\index{Python-Umgebung}

Der einfachste Weg zu einer funktionierenden Umgebung:

\begin{schritte}
  \item Installieren Sie einmal ein Python (3.10 oder neuer) für Ihr System
        (Kapitel~\ref{kap:installation}).
  \item Tragen Sie dieses Python unter \gui{Python interpreter:} ein (oder lassen Sie den
        gefundenen Interpreter stehen).
  \item Klicken Sie auf \gui{Create environment…}. Ein Dialog \gui{New Python environment
        (folder)} fragt nach einem \emph{neuen} Ordner; vorgeschlagen wird
        \datei{python-env} im Datenordner der Anwendung. Existiert der Ordner schon, bricht der
        Vorgang mit \code{… exists already – choose a new folder.} ab.
  \item Die Anwendung legt mit dem Interpreter im Feld eine venv an und installiert darin mit
        \code{pip} \code{numpy} und alle Pakete, die die geladenen Plugins verlangen
        (\code{raytatouille}, \code{rtt-fusion}, \code{photonpower-fieldtracing},
        \code{fouriergrating}, \code{scipy}, \code{pathtracer}). Die Ausgaben erscheinen
        laufend im Textfeld. Das dauert je nach Internetverbindung einige Minuten.
  \item Am Ende steht \code{Done – the new interpreter is set above (applied with OK).} Der
        Interpreter der neuen Umgebung steht jetzt im Feld.
  \item Klicken Sie auf \gui{OK}. Die Plugins werden mit der neuen Umgebung neu geladen.
  \item Zur Kontrolle: \menu{Plugins,Settings…} $\to$ \gui{Check environment} sollte nur
        \checkmark{} zeigen.
\end{schritte}

Schlägt die Installation der Pakete fehl (z.\,B.\ ohne Internetverbindung), steht dort
\begin{quote}\raggedright
\code{Installing the packages failed – see above; Check environment shows what is missing.}
\end{quote}
Der neue Interpreter wird trotzdem eingetragen; holen Sie die fehlenden Pakete später mit dem
von \gui{Check environment} angezeigten Befehl nach (Anhang~\ref{sec:programme-python}).

\subsection{Der Entwicklermodus}\label{sec:plugin-entwicklermodus}\index{Entwicklermodus}

Normalerweise lädt die Anwendung nur Plugins, die mitgeliefert oder vom Herausgeber signiert
sind. Der \emph{Entwicklermodus} lädt zusätzlich unsignierte Plugins und solche mit ungültiger
Signatur – gedacht für Plugin-Entwickler, die ihr eigenes Plugin ausprobieren. Sie schalten ihn
mit dem Kontrollkästchen \gui{Developer mode: also load unsigned plugins} ein (Tooltip:
\gui{Unsigned plugins run Python code that nobody has vouched for. Only for plugins you are
developing.}) oder mit der Umgebungsvariablen \code{NFCAD\_DEVELOPER\_MODE=1}; im zweiten Fall
ist das Kästchen ausgegraut und trägt den Zusatz \gui{(on through NFCAD\_DEVELOPER\_MODE)}.

Solange der Entwicklermodus an ist, steht rechts in der Statuszeile das rote Banner
\gui{DEVELOPER MODE – unsigned plugins allowed}; unsignierte Plugins tragen den Zusatz
\gui{[unsigned]} bzw.\ \gui{[signature invalid]}. Einzelheiten:
Anhang~\ref{sec:verwaltung-entwicklermodus}.

% ---------------------------------------------------------------------------
\section{Plugins aus den Seitenleisten nehmen}\label{sec:plugin-seitenleisten}

Nutzen Sie ein Plugin selten, können Sie es aus der Werkzeugspalte und dem Plugin-Dock nehmen,
ohne es zu deinstallieren:

\begin{schritte}
  \item Öffnen Sie \menu{Plugins,Show in the side panels}.
  \item Entfernen Sie das Häkchen vor dem Plugin (Tooltip: \gui{Tool box group and Plugin panel;
        the actions stay in this menu}).
\end{schritte}

Die Werkzeuggruppe des Plugins verschwindet, und das Plugin fehlt in der Auswahlliste des
Plugin-Docks. Das Plugin bleibt aber geladen: Seine Aktionen stehen weiter im Menü
\menu{Plugins}, es liefert weiter Blöcke für den Node-Workflow, und Aktionen an Blöcken bleiben
im Kontextmenü des Node-Editors. Die Einstellung gilt für die Anwendung (nicht je Dokument) und
bleibt über einen Neustart erhalten. Ein erneutes Häkchen holt das Plugin zurück.

\begin{tipp}
Das Plugin \plugin{Workflow} brauchen Sie meist nur über seine beiden Aktionen
\gui{Check workflow} und \gui{Run workflow} und seine Bildeinstellungen. Wenn Sie den
Node-Workflow nicht nutzen, können Sie es hier aus den Seitenleisten nehmen.
\end{tipp}

% ---------------------------------------------------------------------------
\section{Die mitgelieferten Plugins}\label{sec:plugins-mitgeliefert}\index{mitgelieferte Plugins}

Mit \produkt{} werden fünf Plugins installiert. Sie liegen im schreibgeschützten
Programmordner und werden ohne Signaturprüfung geladen, weil sie Teil der Installation sind
(Anhang~\ref{sec:verwaltung-mitgeliefert}). Ihre Python-Pakete installiert
\gui{Create environment…} (Abschnitt~\ref{sec:umgebung-anlegen}).

\begin{referenz}{@{}p{0.2\linewidth}p{0.13\linewidth}p{0.12\linewidth}X@{}}
\toprule
\textbf{Plugin} & \textbf{Kennung} & \textbf{Version} & \textbf{Benötigte Python-Pakete} \\
\midrule
\endhead
\plugin{Raytracer} & \code{raytracer} & 0.1.0 & \code{numpy}, \code{raytatouille>=0.9},
  \code{rtt\_fusion>=0.2} (pip: \code{rtt-fusion}), \code{fieldtracing>=0.2}
  (pip: \code{photonpower-fieldtracing}) \\
\plugin{Workflow} & \code{workflow} & 0.1.0 & \code{numpy}, \code{rtt\_fusion>=0.2},
  \code{raytatouille>=0.9} \\
\plugin{FourierGrating} & \code{fg} & 0.1.0 & \code{numpy}, \code{scipy},
  \code{fouriergrating>=0.18} \\
\plugin{Path Tracer} & \code{pathtracer} & 0.1.0 & \code{numpy}, \code{pathtracer>=0.3} \\
\plugin{Automotive Lighting} & \code{al} & 0.22.0 & \code{numpy}, \code{pathtracer>=0.3} \\
\bottomrule
\end{referenz}

\paragraph{Raytracer.} Verbindet den Strahlverfolger \emph{Raytatouille}. Er lädt optische
Systeme aus \datei{.rtt.json}-Dateien (\gui{Load system}) und zeichnet ihre Flächen als exakte
NURBS-Flächen, verfolgt Strahlen und zeigt sie als Linien (\gui{Trace rays}), berechnet mit
vielen Strahlen ein Kamerabild (\gui{Camera image}), optimiert das System mit live
aktualisierten Flächen (\gui{Optimise}) und speichert es (\gui{Save system}). Außerdem macht er
aus einem System einen Node-Workflow (\gui{Workflow from system}, \gui{New Raytatouille
system}), liefert die Raytatouille-Blöcke \code{rtt.*} und hält im Feld \gui{Glass catalogues}
die Glaskataloge für alle Rechnungen. Kapitel~\ref{kap:raytracer},
Anhang~\ref{kap:plugin-raytracer}.

\paragraph{Workflow.} Die Rechenmaschine des Node-Workflows. Sie prüft den Workflow
(\gui{Check workflow}), plant, in welcher Darstellung das Licht über jede Verbindung läuft, und
rechnet, was die Anzeigeblöcke zeigen sollen (\gui{Run workflow}). Ihr Panel enthält die
Bildeinstellungen (Strahlen, Pixel, Bildfenster, Gauß-Beleuchtung, Feldgenauigkeit, Wellenoptik
nahe einem Fokus). Sie kennt selbst keinen Rechenkern, sondern lässt die Blöcke von den Plugins
rechnen, die sie liefern. Kapitel~\ref{kap:workflow}, Anhang~\ref{kap:workflow-bloecke}.

\paragraph{FourierGrating.} Beugung an periodischen Strukturen (RCWA, differentielle Methode).
Im Panel beschreiben Sie Struktur, Materialien, Schichten und Einfall; \gui{Solve} zeigt die
Effizienz jeder Ordnung als Tabelle und Balken, \gui{Scan} Effizienzen über Wellenlänge,
Winkel oder Polarisation als Kurven, \gui{Near field} das Feld im Gitter als Bilder,
\gui{Optimise} findet Schichtdicken, Stützstellen oder Perioden für eine Zielfunktion. Dazu
kommen das Laden und Speichern von Konfigurationen, Schichten aus CAD-Körpern und die Blöcke
\code{fg.grating} und \code{fg.stack} für den Node-Workflow. Kapitel~\ref{kap:fouriergrating},
Anhang~\ref{kap:plugin-fouriergrating}.

\paragraph{Path Tracer.} Lichttransport direkt auf dem CAD-Modell: Objekte werden zu
Lichtquellen, die übrigen bekommen Materialien (diffus, Spiegel, glänzend, Glas, Metall,
Mattglas, Absorber, Detektor). \gui{Irradiance} legt die Bestrahlungsstärke als Textur auf die
Flächen, \gui{Camera image} berechnet das Bild einer Kamera – mit Light Tracing, Path Tracing,
bidirektionalem Path Tracing oder Photon Mapping. Mit dem Paket \code{pathtracer} rechnet es in
C++ auf allen Prozessorkernen. Kapitel~\ref{kap:pathtracer},
Anhang~\ref{kap:plugin-pathtracer}.

\paragraph{Automotive Lighting.} Scheinwerfer und Signalleuchten für Glühlampen: der
facettierte Reflektor eines Klarglas-Scheinwerfers, dessen Facetten die Hell-Dunkel-Grenze
bilden, und die Streuscheibe vor einem Paraboloid. Das Plugin erzeugt Reflektor und Lampe als
NURBS-Flächen, zeigt Wendelbilder und Lichtverteilung mit den Prüfpunkten nach UN/ECE R112 und
optimiert Streuscheiben; die Lichtverteilung lässt sich auch mit dem Rechenkern des Path Tracers
strahlverfolgen. Kapitel~\ref{kap:automotive}, Anhang~\ref{kap:plugin-automotive}.

% ---------------------------------------------------------------------------
\section{Typische Probleme}\label{sec:plugins-probleme}

\begin{referenz}{@{}p{0.36\linewidth}X@{}}
\toprule
\textbf{Beobachtung} & \textbf{Ursache und Abhilfe} \\
\midrule
\endhead
Im Plugin-Dock steht \gui{not running – see the plugin log}. &
  Das Plugin ließ sich nicht starten. Öffnen Sie \menu{View,Plugin log}: Meist fehlt ein
  Paket (\code{ModuleNotFoundError}) oder der Interpreter ist falsch. Prüfen Sie mit
  \gui{Check environment} und installieren Sie das Fehlende, dann
  \menu{Plugins,Reload plugins}. \\
Ein Klick auf eine Aktion bewirkt nichts. &
  Ein Feld im Formular ist rot markiert (Zahlenliste, Tabelle, JSON). Korrigieren Sie es. \\
\gui{A plugin job is still running} &
  Es läuft noch ein Job, eventuell eines anderen Plugins. Warten Sie oder brechen Sie mit
  \keys{Esc} ab. \\
Ein Meldungsfenster mit einem Text des Plugins erscheint. &
  Ein Fehler, den Sie beheben können (fehlende Datei, falsche Eingabe); der Text sagt, was zu
  tun ist. \\
\textit{Plugin} \code{failed: …} und das Log öffnet sich. &
  Ein Programmfehler oder ein Problem des Rechenkerns; der Traceback steht im Log. Hängen Sie
  ihn an eine Fehlermeldung an. \\
Das Plugin fehlt in Dock und Werkzeugspalte, steht aber im Menü. &
  Es ist unter \menu{Plugins,Show in the side panels} abgewählt. \\
Ein eigenes Plugin wird nicht geladen; unter \gui{Problems:} steht \code{not signed}. &
  Nur signierte Plugins werden geladen; zum Entwickeln den Entwicklermodus einschalten. \\
Der Abbruch dauert einige Sekunden. &
  Das Plugin bricht erst an der nächsten geeigneten Stelle ab; nach 5~s wird sein Prozess
  beendet. \\
\bottomrule
\end{referenz}

Weitere Fehlerbilder sammelt Anhang~\ref{kap:fehlersuche}.
